\section{Source bounds and retained routes}
\label{sec:global-construction}

We now use the finite-host construction of Zhao--Booth-Clibborn--Schwahn
\cite[Sections~4--5]{ZBS}, built on Lu--Wang's retained-set descent
\cite[Theorem~2]{LW}. Proposition~\ref{prop:local-alternative} supplies
its improved local premise. The closure, route composition, and joint
induction retain the earlier structure.

Call a positive pair $(A,B)$ a \emph{source} if one constant $C\ge1$
works for all positive integers $k,\ell$ in
\[
 R(k,\ell)\le C e^{Ak+B\ell}.
\]
\leannote{\texttt{SourceAdmissible},
\leanfile{S1_Definitions/L2_AdmissibleBounds.lean\#L20}{L2_AdmissibleBounds.lean}}
A \emph{profile} $P=(p_P,I_P,L_P)$ retains a density hypothesis.
Here $0<p_P<1$, $I_P\subset(0,\infty)$ is a compact interval, and
$L_P$ is positive on $I_P$. Write $\rho=\ell/k$ for the target ratio
and $d_R(G)$ for the red edge density.
The profile asserts, uniformly in $k,\ell$ and the graph $G$,
\begin{equation}\label{eq:profile}
 d_R(G)\ge p_P,\quad \rho\in I_P,\quad
 |V(G)|\ge C e^{kL_P(\rho)}
 \quad\Longrightarrow\quad
 G\text{ contains a red }K_k\text{ or a blue }K_\ell.
\end{equation}
\leannote{\texttt{BankProfile.Assert},
\leanfile{Support/Published/Definitions/Def_DiagRamsey_JointBank.lean\#L46}{Def_DiagRamsey_JointBank.lean}}
The same $C$ must work throughout the ratio interval.

\subsection{Closure and route outputs}

A \emph{control} $\gamma$ consists of parameters satisfying
Proposition~\ref{prop:local-alternative}, a terminal source
$(A_\sigma,B_\sigma)$, and numbers
\[
 0<\xi<\min\{x,e^{-A_\sigma}\},\qquad
 0<y_0<e^{-B_\sigma},\qquad 0<\nu<\mu.
\]
\leannote{\texttt{ComplementControlValid},
\leanfile{S6_GlobalConstruction/L1_ProfileSupport.lean\#L26}{L1_ProfileSupport.lean}}
Set $a_\gamma=-\log\xi$, $b_\gamma=-\log y_0$ and
$c_\gamma=-\log\nu$. These are positive. At this stage the terminal
source is a proposed bound, to be proved together with the profiles.

The proof of the finite-host closure theorem~\cite[Theorem~4.2 and
Section~4.3]{ZBS} applies with Proposition~\ref{prop:local-alternative}
in place of its local alternative. It requires the source only on
proper subsets of the current host: for every $a\ge1$,
a proper subset of size at least $C e^{A_\sigma a+B_\sigma\ell}$
must contain a red $K_a$ or a blue $K_\ell$.
Under this premise, there is $L_0$ depending only on the control such
that, for $\ell\ge L_0$, $k,b\ge1$, and disjoint nonempty $X,Y$ of
red cross-density at least $p_\gamma$,
\begin{equation}\label{eq:closure-volume}
 |X|^w|Y|\ge C^{w+1}\xi^{-k}y_0^{-\ell}\nu^{-wb}
 \quad\Longrightarrow\quad
 \begin{cases}
 \text{a red }K_k\text{ in }X\cup Y,\text{ or}\\
 \text{a blue }K_b\text{ in }X,\text{ or}\\
 \text{a blue }K_\ell\text{ in }Y.
 \end{cases}
\end{equation}
\leannote{\texttt{finite\_host\_closure},
\leanfile{S6_GlobalConstruction/L3_RouteComposition.lean\#L34}{L3_RouteComposition.lean}}
The proof fixes the moment parameter and initializes the potential at a
shift just below $p_\gamma$, and repeatedly applies maximal restriction
and the local alternative. A successful test reduces a clique target.
The case of many large blue neighbourhoods is handled by a blue book
of logarithmic size. The strict inequalities for $\xi,y_0,\nu$ ensure
the required bounds despite the initial restriction, blue-book construction,
and integer rounding. The stopping rule is invoked only on proper
subhosts. Thus the proof of~\cite[Section~4.3]{ZBS} applies with our local
alternative; the closure induction originates in~\cite[Lemma~A.1]{LW}.

A balanced cut preserving red density gives the \emph{leaf output}
\[
 L_\gamma(\rho)=
 \frac{a_\gamma+\rho(b_\gamma+wc_\gamma)}{1+w}.
\]
\leannote{\texttt{SharpControl.leafLine},
\leanfile{Support/Published/Definitions/Def_DiagRamsey_SourceBank.lean\#L52}{Def_DiagRamsey_SourceBank.lean}}
Indeed, both parts have size at least one third of the host.
At host size $3Ce^{kL_\gamma(\rho)}$, the volume condition
\eqref{eq:closure-volume} holds with $b=\ell$.
Any fixed positive margin above $L_\gamma$ absorbs the factor $3$.

A \emph{retained route} improves this output by calling previously
specified profiles along the descent. Choose knots
$0<\theta=u_0<\cdots<u_M=\omega$, a starting height $g_0>0$, and
slopes $d_j$, and put
\[
 g(\rho)=g_0+\sum_{j=0}^{M-1}
 d_j\bigl(\min\{\rho,u_{j+1}\}-u_j\bigr)_+.
\]
\leannote{\texttt{RoutePiece.g},
\leanfile{Support/Published/Definitions/Def_DiagRamsey_JointBank.lean\#L68}{Def_DiagRamsey_JointBank.lean}}
On $[u_j,u_{j+1}]$, call a child $P_j$ whose domain contains that cell,
with $L_j(\rho)<g(\rho)$ and $d_j>-\log(1-p_j)$.
If a retained graph has density at least $p_j$, the child gives the
desired clique. Otherwise a vertex has a blue neighbourhood of size
at least $(1-p_j)(|Z|-1)$. Since $e^{-d_j}<1-p_j$, this is at least
$e^{-d_j}|Z|$ for sufficiently large $|Z|$. The neighbourhood therefore
satisfies the size bound for the blue target reduced by one.
Bounded corrections at the knots account for integer targets.

The route output is
\begin{equation}\label{eq:route-output}
 V_\gamma(\rho)=\max\left\{
 g(\rho),\,
 \frac{a_\gamma+\rho b_\gamma+
 w\bigl(\theta c_\gamma+g(\rho)-g_0\bigr)}{1+w}\right\}.
\end{equation}
\leannote{\texttt{RoutePiece.value},
\leanfile{Support/Published/Definitions/Def_DiagRamsey_JointBank.lean\#L72}{Def_DiagRamsey_JointBank.lean}}
The first term enforces the size condition at the start of the route;
the second enforces the closure condition at its terminal ratio $\theta$.
A profile of density strictly greater
than $p_\gamma$ and height exceeding $V_\gamma$ by a fixed positive
margin is therefore valid on its covered interval. The margin absorbs
the initial restriction and the bounded knot corrections.
Every child is called on a subset of $X$, disjoint from the nonempty
set $Y$, hence on a proper subhost. This is the route argument
of~\cite[Definition~5.4 and Lemma~5.8]{ZBS}; its implementation is
\leanfile{S6_GlobalConstruction/L3_RouteComposition.lean\#L445}{L3_RouteComposition.lean}.

\subsection{Proving sources and profiles together}

Source bounds and profiles may depend on one another. The dependence
is resolved by induction on host order, as in~\cite[Theorem~5.7]{ZBS}.
Take finite families of proposed positive source pairs and profiles.
Call profiles obtained from leaves, retained routes, and proper-subhost
complementary pairs \emph{direct}. Call those obtained from direct
profiles by density transfer, complementary profile pairs, and monotonicity
in the blue clique target \emph{combined}. The proper-subhost complementary
pairs and target-monotonicity rules extend the list
in~\cite[Definitions~5.4--5.5]{ZBS}; their precise hypotheses are in
\leanfile{Support/Published/Definitions/Def_DiagRamsey_JointBank.lean\#L78}{Def_DiagRamsey_JointBank.lean}.
Each profile domain must be covered by valid pieces, with the required
output margins.

Let $F(\rho)$ be the symmetric Lu--Wang baseline exponent, whose
uniform source estimate permits an arbitrarily small positive error proportional to
$k+\ell$~\cite[Lemma~B.1]{LW}. For each proposed source $(A,B)$,
choose $\epsilon>0$. At every $\rho>0$, require either a baseline
witness or two profiles with $p_R+p_B\le1$, valid at $\rho$ and
$1/\rho$, respectively, such that
\[
 F(\rho)\le A+B\rho-\epsilon(1+\rho)
 \quad\text{or}\quad
 \max\{L_R(\rho),\rho L_B(1/\rho)\}
       \le A+B\rho-\epsilon(1+\rho).
\]
\leannote{\texttt{SourceCell.Valid},
\leanfile{Support/Published/Definitions/Def_DiagRamsey_JointBank.lean\#L128}{Def_DiagRamsey_JointBank.lean}}
The density condition ensures that one of the two profile assertions
applies in every colouring.

Choose one constant $C$ to cover all small-target cases and the finitely
many baseline estimates. This is possible because the profile ratio
domains are bounded, and, when one source target is bounded, the
elementary binomial Ramsey bound grows polynomially in the other;
positive source exponents absorb this polynomial.
Now induct on the host order $N$. On every proper subhost, the induction
hypothesis supplies all proposed sources and profiles. First prove the
direct profiles at order $N$ by their covering pieces. Next prove the
combined profiles using those direct profiles at the same order.
Finally prove the sources, using either the baseline estimate or the
two profiles in the source witness. Thus each recursive call decreases
the host order, while same-host calls follow the order
direct profiles, combined profiles, sources.

This proves all assertions with one common prefactor. Rounding host
thresholds upward enlarges it by at most one and gives the stated
Ramsey source bounds. See \texttt{joint\_profiles\_of\_pieces} in
\leanfile{S6_GlobalConstruction/L2_JointInduction.lean\#L28}{L2_JointInduction.lean}.
The next section supplies the finite data and strict inequalities for
the construction used here.
